\documentclass[12pt]{article}
\usepackage[margin=1in]{geometry}
\usepackage{amsmath,amssymb,amsfonts}
\usepackage{graphicx}
\usepackage{booktabs}
\usepackage{hyperref}
\usepackage{authblk}
\usepackage{setspace}
\onehalfspacing

\newcommand{\si}[1]{\mathrm{#1}}

\title{Environment Structure, Not Just Team Size: A Case Study in the\\
Generalization Limits of Permutation-Invariant Coverage Policies}
\author[1]{Harsh Pandhe}
\affil[1]{Independent Research \\ \texttt{harshpandhehome@gmail.com}}
\date{}

\begin{document}
\maketitle

\begin{abstract}
\noindent
Permutation-invariant neural architectures (deep sets, graph neural networks)
are well established as a route to multi-robot policies that generalize
across team size without retraining. We highlight a second, less examined
axis of generalization: environment \emph{structure}. Using an identical
navigation stack -- an A*-augmented frontier-exploration controller behind a
permutation-invariant neighbor encoder -- deployed unchanged into two
real-world-derived Gazebo environments, we show that reachable area-coverage
ceiling differs sharply by environment structure alone: $90.4\%$ in a densely
furnished cafe environment versus $99.0\%$--$100.0\%$ in an open warehouse
environment containing only isolated pillar obstacles, of comparable or
larger physical area. Both environments were independently verified via
offline mesh/geometry analysis before deployment, so the difference is
attributable to environment topology (dense, coverage-blocking clutter versus
sparse, routable obstacles) rather than measurement artifacts. We argue that
team-size invariance and environment-structure invariance are separate
properties that do not follow from each other, and that reporting a single
coverage number without characterizing the environment's obstacle topology
understates how much of a policy's apparent performance is actually
attributable to the environment it was tested in.
\end{abstract}

\noindent\textbf{Keywords:} multi-robot coverage, permutation-invariant
neural networks, graph neural networks, generalization, swarm robotics

\section{Introduction}
A permutation-invariant policy -- one whose neighbor-observation processing
is invariant to the number and ordering of neighbors, typically via a shared
per-neighbor encoder followed by a symmetric pooling
operation~\cite{zaheer2017deepsets} -- is now a standard tool for building
multi-robot policies that scale to team sizes unseen during
training~\cite{multientitypinn2024,localcanon2025,gnnrealworld2021}. This is
a genuinely useful and well-demonstrated property. Our concern in this paper
is a different, easily conflated one: does scaling to a new \emph{team size}
with a fixed architecture say anything about scaling to a new
\emph{environment}? We argue it does not, by construction -- team-size
invariance is a property of the neighbor-processing module, while
environment generalization additionally depends on the exploration/planning
components consuming the environment's structure, which permutation
invariance does not touch.

We make this concrete with a controlled case study: the same navigation
stack, deployed unchanged into two physically distinct, independently
verified environments -- a furnished cafe and an open warehouse -- and show
that the ceiling on achievable area coverage differs by roughly $9$--$20$
percentage points between them, driven entirely by environment topology.

\section{Related Work}
Zaheer et al.~\cite{zaheer2017deepsets} establish the deep-sets construction
--- shared per-element encoding followed by a symmetric pooling operation ---
as the canonical route to permutation-invariant function approximation, which
underlies most permutation-invariant multi-robot policies including ours. An
et al.~\cite{multientitypinn2024} apply permutation-invariant networks to
multi-entity robotic control via a PointNet-style architecture. Wang et
al.~\cite{localcanon2025} combine agent-centric canonicalization with
role-aware graph encoders for equivariant, sample-efficient, and
generalizable swarm control, validated in a sim-to-real pursuit-evasion
setting. Blumenkamp et al.~\cite{gnnrealworld2021} deploy decentralized
GNN-based policies on a real multi-robot system communicating over ad-hoc
networks, demonstrating sim-to-real transfer of a GNN policy architecture.
Gosrich et al.~\cite{gosrich2022coverage} apply GNNs directly to multi-robot
coverage control, the same task family we study, and report improved
coverage quality from inter-robot communication relative to purely local
information -- complementary to our finding, since their generalization axis
is communication topology and team size, not physical environment structure.
Muthusamy et al.~\cite{gnngeneralizability2024} study GNN generalizability
for decentralized unlabeled motion planning, the closest prior work to our
concern, but frame generalization primarily in terms of team size and graph
connectivity rather than environment obstacle topology. To our knowledge, no
prior work in this family holds team size and architecture fixed while
varying environment structure as the sole independent variable, which is
precisely the controlled comparison we report.

\section{Policy Architecture}
\label{sec:arch}
The actor separates ego features $\mathbf{e}\in\mathbb{R}^{28}$ (LiDAR, goal,
velocity) from a variable-size neighbor set. Each neighbor's relative
$(\text{distance}, \text{angle})$ pair is embedded by a shared two-layer MLP
$\phi:\mathbb{R}^2\to\mathbb{R}^{16}$; padded (absent) neighbors are masked
out and the valid embeddings mean-pooled:
\begin{equation}
\mathbf{g} = \frac{1}{\max(1,|\mathcal{V}|)}\sum_{j\in\mathcal{V}} \phi(n_j),
\qquad \mathcal{V}=\{j: d_j < 9.0\,\si{m}\},
\end{equation}
which is invariant to neighbor count and ordering by construction, and was
verified to run unchanged (no dimension mismatch, correct output shapes) at
$N=5$ and $N=10$ neighbors against a network trained assuming $N=3$. This
module is architecturally identical across both environments studied below;
what differs between environments is not this module but the frontier
target-selection and A* routing logic each robot uses to decide
\emph{where} to go, which necessarily depends on the map it observes.

\section{Case Study Design}
\label{sec:design}
We deploy an identical three-robot navigation stack -- frontier-based target
selection (closest unvisited grid cell) augmented with A* routing (activated
only when the direct line to the target is obstructed, with a one-cell
obstacle-inflation safety margin) -- into two Gazebo Sim environments.

\textbf{Cafe environment.} A furnished cafe model (Gazebo Fuel,
\texttt{openrobotics/models/cafe}) with tables, counters, and other clutter
distributed through the interior. Coverage grid: $40\times20$ cells over a
$\sim17\times8\,$m usable region ($\sim136\,\si{m}^2$).

\textbf{Warehouse environment.} An open warehouse model (Gazebo Fuel,
\texttt{openrobotics/models/Warehouse}), independently verified offline
before deployment by downloading its collision mesh (STL) and computing
occupancy at LiDAR sensing height ($10$--$30\,\si{cm}$): a $30\times50\,$m
overall footprint that is $95$--$100\%$ open floor at that height, with
obstacles limited to outer walls and six isolated support pillars arranged in
a $2\times3$ grid. We deployed into an $18\times12\,$m region ($216\,\si{m}^2$
-- larger than the cafe region) deliberately including one pillar pair (at
world coordinates $x\approx\pm7.25\,$m, $y\approx0\,$m), so the comparison
includes genuine obstacle discovery and avoidance rather than an artificially
obstacle-free room. Coverage grid: $45\times30$ cells, scaled to keep cell
size ($0.4\,$m) comparable to the cafe grid and above the navigation stack's
$0.2\,$m target-arrival tolerance -- reusing a fixed grid resolution across
differently-sized environments was found in preliminary testing to produce
cells smaller than this tolerance, causing a target cell to be approached but
never entered, and therefore never marked visited, a deadlock that froze the
swarm at $3$--$18\%$ coverage; this is itself a second, distinct
environment-structure sensitivity (grid discretization relative to physical
scale) that a fixed architecture does not resolve automatically.

Both environments use the same per-step reward-shaping-free evaluation
protocol: a single long episode (continuous re-targeting on goal arrival, no
episode-ending termination on collision) run to $1200$--$1600$ steps, with
area-coverage rate (ACR), overlap redundancy, swarm path length, and a
collision-step count (steps where any robot's minimum LiDAR range is below
$0.20\,$m) recorded throughout.

\section{Results}
\begin{table}[h]
\centering
\caption{Coverage ceiling by environment, identical navigation stack.}
\label{tab:results}
\begin{tabular}{@{}lccccc@{}}
\toprule
Environment & Area ($\si{m}^2$) & Steps & ACR (\%) & Redund.\ & Coll.-steps \\
\midrule
Cafe (reactive only)        & 136 & 1200 & 87.5 & --    & 463  \\
Cafe (+ A* routing)         & 136 & 1200 & 90.4 & --    & 1824 \\
Warehouse (2 real pillars)  & 216 & 1600 & 99.0 & 10.11 & 607  \\
\bottomrule
\end{tabular}
\end{table}

The cafe environment's coverage ceiling under a purely reactive controller
(no path planning around obstacles) is $87.5\%$, and does not improve with
additional run length -- the remaining $\sim$12\% consists of cells the
controller can see but cannot reach in a straight line, and reactive-only
avoidance cannot route around the intervening furniture. Adding A* routing
recovers some of this gap ($87.5\%\to90.4\%$) at the cost of a substantially
higher collision-step count, since the swarm now pushes into tighter,
harder-to-reach corners near obstacles for the remaining coverage. The
warehouse environment, despite covering a larger physical area and containing
genuine pillar obstacles the swarm must discover and route around, reaches
$99.0\%$ under the identical A*-augmented stack -- because its obstacles are
sparse and isolated rather than densely arranged, essentially no floor area
is rendered permanently unreachable by the environment's topology.

\section{Discussion}
The permutation-invariant neighbor-encoding module described in
Section~\ref{sec:arch} is unchanged across both environments and both A*
configurations in Table~\ref{tab:results} -- the $9$--$20$ percentage-point
coverage gap is not attributable to it. It is attributable entirely to
environment topology interacting with the (non-invariant) exploration and
routing logic: a environment whose obstacles create disconnected or
hard-to-reach pockets imposes a coverage ceiling that no amount of
neighbor-encoder generalization addresses, because the neighbor encoder never
sees the static environment map at all -- it only processes other robots'
relative states. This is the sense in which team-size invariance and
environment-structure invariance are independent properties: an architecture
can be perfectly invariant to how many teammates are present while still
being highly sensitive to where the walls are.

A second, narrower finding reinforces the same point at the level of
implementation detail rather than architecture: the coverage-grid deadlock
described in Section~\ref{sec:design} shows that even a supporting
data-structure choice (grid cell size) that was well-tuned for one
environment's physical scale produced a hard failure (not merely reduced
performance) when reused unchanged in a differently-scaled environment. Grid
resolution is not part of the neural architecture at all, but it is exactly
the kind of environment-coupled implementation detail that a "team-size
generalization" evaluation would never surface, since it depends on absolute
physical scale rather than agent count.

We see two practical implications. First, evaluations of permutation-invariant
multi-robot policies that vary team size while holding environment fixed (the
common protocol in the cited related work) are measuring a real but partial
form of generalization; a companion evaluation varying environment structure
while holding team size fixed, as done here, measures an orthogonal
dimension that is at least as relevant to real deployment, where the
environment is rarely a design choice. Second, for coverage-specific tasks,
progress on the coverage ceiling itself (Table~\ref{tab:results}, cafe
$87.5\%\to90.4\%$) came entirely from the non-invariant planning component
(A* routing), not from the invariant one -- suggesting that closing the
remaining gap toward $100\%$ in structurally harder environments is a
planning-and-representation problem for the \emph{map}, not a scaling
problem for the \emph{team}.

\section{Limitations and Future Work}
This is a two-environment case study, not a systematic sweep; we do not
claim the specific $9$--$20$ point gap generalizes quantitatively to other
environment pairs, only that a gap of this kind is real, reproducible, and
attributable to topology rather than team size or architecture. Both
environments are simulated; the natural extension is sim-to-real hardware
transfer, following the real-world GNN deployment pattern demonstrated by
Blumenkamp et al.~\cite{gnnrealworld2021}, to test whether the same
structure-dependent ceiling reproduces physically, and a broader sweep across
environment topologies (varying obstacle density and connectivity
independently) to characterize the relationship between topology and
coverage ceiling more precisely than the two-point comparison given here.

\section{Conclusion}
Holding a permutation-invariant multi-robot policy's architecture and team
size fixed while varying the deployment environment reveals a
$9$--$20$-percentage-point swing in achievable area coverage, attributable to
environment topology rather than the neighbor-encoding module responsible
for team-size generalization. We argue this constitutes a second,
independent axis of generalization for multi-robot coverage policies that
existing team-size-focused evaluations do not surface, and that reporting a
coverage number without characterizing the deployment environment's obstacle
topology risks attributing environment-driven performance differences to the
policy itself.

\bibliographystyle{plain}
\bibliography{topic3_references}

\end{document}
